\documentclass{article}
\usepackage[margin=1in]{geometry} % margins
\usepackage[justification=centering]{caption}
%\usepackage{calrsfs}
\usepackage{amsmath}
\usepackage{amsfonts}
\usepackage{amssymb}
\usepackage{parskip}
\usepackage{graphicx}
\usepackage{enumerate}
\usepackage{enumitem}
\usepackage{amsthm}
\usepackage{float}
\usepackage{tikz-cd}
\usepackage{ytableau}
\usepackage{dsfont}

\renewcommand{\implies}{\Rightarrow}
\newcommand{\longimplies}{\Longrightarrow}
\newcommand{\N}{\mathbb{N}}
\newcommand{\Z}{\mathbb{Z}}
\newcommand{\Q}{\mathbb{Q}}
\newcommand{\R}{\mathbb{R}}
\newcommand{\C}{\mathbb{C}}
\newcommand{\D}{\mathbb{D}}
\renewcommand{\H}{\mathbb{H}}
\newcommand{\F}{\mathbb{F}}
\newcommand{\K}{\mathbb{K}}
\newcommand{\Lf}{\mathbb{L}}
\newcommand{\Sp}{\mathbb{S}}
\newcommand{\cl}[1]{\overline{#1}}
\renewcommand{\emptyset}{\varnothing}
\newcommand{\norm}[1]{\left\|#1\right\|}
\newcommand{\maxnorm}[1]{\left\|#1\right\|_{\text{max}}}
\newcommand{\opnorm}[1]{\left\|#1\right\|_{\mathrm{op}}}
\newcommand{\supnorm}[1]{\left\|#1\right\|_{\infty}}
\newcommand{\id}{\operatorname{id}}
\newcommand{\sspan}{\operatorname{span}}
\newcommand{\tr}{\operatorname{tr}}
\newcommand{\ch}{\operatorname{char}}
\newcommand{\poly}{\mathrm{P}}
\newcommand{\mtx}{\mathcal{M}}
\newcommand{\seq}{\mathbf{Seq}}
\newcommand{\vv}[1]{\mathbf{#1}}
\newcommand{\dirsum}{\oplus}
\newcommand{\rank}{\operatorname{rank}}
\newcommand{\im}{\operatorname{im}}

\newcommand{\cj}[1]{\overline{#1}}
\newcommand{\inprod}[2]{\left\langle#1,#2\right\rangle}
\renewcommand{\Re}{\operatorname{Re}}
\renewcommand{\Im}{\operatorname{Im}}

\newcommand{\res}{\operatorname{res}}

\newcommand{\floor}[1]{\left\lfloor#1\right\rfloor}

% Theorem environments
\usepackage{amsthm}

\theoremstyle{definition} % The "plain" style italicizes all body text.
	\newtheorem{thm}{Theorem}
		\numberwithin{thm}{section} % Theorem numbers are determined by section.
	\newtheorem{lemma}[thm]{Lemma}
	\newtheorem{prop}[thm]{Proposition}
	\newtheorem{cor}[thm]{Corollary}
	\newtheorem{conj}[thm]{Conjecture}

\theoremstyle{definition} % The "definition" style does not italicize body text..
    \newtheorem{defn}[thm]{Definition}
	\newtheorem{example}[thm]{Example}
    \newtheorem{exercise}[thm]{Exercise}
    \newtheorem{question}[thm]{Question}

\usepackage{mathtools}

\usetikzlibrary{decorations.markings,positioning,decorations.text}

\def\xr{5}
\def\yr{5}

%\usepackage{titlesec}
%\titleformat{\section}[hang]{\normalfont\fontsize{14}{0}\bfseries}{Problem 8.\thesection\ }{0pt}{}

\begin{document}

\section{Worksheet}

\begin{question}
Let $E \in \mathcal{E}$ satisfy $m(E) > 0$. Show that there is a bounded set $A \subseteq E$ such that $m(A) > 0$.
\end{question}
\begin{proof}
Consider the increasing collection ${\{A_n := (-n, n) \cap E\}}_{n = 1}^\infty$ of measurable sets. By continuity of measure, we have
\[ \lim_{n \to \infty} m(A_n) = m\left(\bigcup_{n = 1}^\infty A_n\right) = m(E) > 0 \]
Since the limit of $m(A_n) > 0$, we must have that $m(A_n) > 0$ for at least one $n$. Hence $A := (-n, n) \cap E \subseteq E$ is a bounded set with $m(A) > 0$.
\end{proof}

\begin{question}
Show that every open set $O \subset \mathbb{R}$ is a countable union of disjoint intervals.
\end{question}
\begin{proof}
For each $x \in O$, there is some $\varepsilon > 0$ such that $(x - \varepsilon_x, x + \varepsilon_x) \subseteq O$. Note that if $\delta > 0$ is such that for all $0 < \varepsilon < \delta$, $(x - \varepsilon, x + \varepsilon) \subseteq O$, then $(x - \delta, x + \delta) = \bigcup_{0 < \varepsilon < \delta} (x - \varepsilon, x + \varepsilon) \subseteq O$. We thus let $\delta_x = \sup\{\varepsilon > 0: (x - \varepsilon, x + \varepsilon) \subseteq O\} > 0$, which is finite since $O \neq \R$ and satisfies $U_x := (x - \delta_x, x + \delta_x) \subseteq O$.

By construction of the $\delta_x$'s, we have $\bigcup_{x \in O \cap \Q} U_x$, a countable union of disjoint intervals, is a subset of $O$. For the reverse equality, for any point $y \in O$, we take $\varepsilon > 0$ for which $(y - 2\varepsilon, y + 2\varepsilon) \subseteq O$, and we take some rational point $x \in (y - \varepsilon, y + \varepsilon) \subseteq O$. Note that then $y \in (x - \varepsilon, x + \varepsilon) \subseteq (y - 2\varepsilon, y + 2\varepsilon) \subseteq O$, so by construction, $y \in (x - \varepsilon, x + \varepsilon) \subseteq U_x$. Therefore, $O \subseteq \bigcup_{x \in O \cap \Q} U_x$, which proves that $O$ is a countable union of disjoint open intervals.
\end{proof}

\begin{question}
Let $E \in \mathcal{E}$ be Lebesgue measurable with $m(E) > 0$. Let $0 < \alpha < 1$. Show that there is a bounded interval $I$ so that
\[ m(E \cap I) \geq \alpha m(I) \]
\end{question}
\begin{proof}
For any $A \subseteq E$, we have $m(E \cap I) \geq m(A \cap I)$ for any bounded interval $I$. By question 1, there is a bounded set $A$ with positive measure, so we may assume that $E$ is already bounded. Take $M > 0$ such that $E \subseteq (-M, M)$. By way of contradiction, suppose that $m(E \cap I) < \alpha m(I)$ for all bounded intervals $I$. For a fixed $\varepsilon > 0$, take open intervals ${\{I_k\}}_{k = 1}^\infty$ that cover $E$ and $\sum_{k = 1}^\infty \ell(I_k) < m(E) + \varepsilon$. We then have 
\[ m(E) \leq \sum_{k = 1}^\infty m(E \cap I_k) < \alpha \sum_{k = 1}^\infty \ell(I_k) < \alpha m(E) + \alpha\varepsilon \]
Letting $\varepsilon \to 0$, we see that $m(E) \leq \alpha m(E)$, which is impossible since $m(E) \in (0, \infty)$ (recall that we took $E$ to be bounded) and $\alpha \in (0, 1)$. Therefore, we must have some bounded interval $I$ for which $m(E \cap I) \geq \alpha m(I)$.
\end{proof}
    
\begin{question}
Let $E \in \mathcal{E}$ be Lebesgue measurable with $m(E) > 0$. Prove that the set $A = E - E := \{x - y : x \in E, y \in E\}$ contains an open interval $J$ such that $0 \in J$.
\end{question}
\begin{proof}
By question 3, we take a bounded interval $I$ such that $m(E \cap I) \geq \frac{3}{4}m(I)$. If we suppose that the result does not hold, then we can take a sequence $x_n \to 0$ such that $x_n \notin E - E$. We take $n$ large enough such that $|x_n| < \frac{1}{4}m(I)$ and note that $x_n + E \subseteq E^c$. We then get the following:
\begin{align*}
    m(I)
    &= m(E \cap I) + m(I \cap E^c) \\
    &\geq m(E \cap I) + m(I \cap (x_n + E)) \\
    &= m(E \cap I) + m((I - x_n) \cap E) \\
    &\geq 2m(E \cap I) - m(((I + x_n) \setminus (I - x_n)) \cap E) \\
    &\geq 2m(E \cap I) - m((I + x_n) \setminus (I - x_n)) \\
    &\geq 2m(E \cap I) - 2|x_n| \\
    &> \frac{3}{2}m(I) - \frac{1}{2}m(I) \\
    &= m(I)
\end{align*}
which is a contradiction.
\end{proof}

\begin{question}
Let $f : \mathbb{R} \to \mathbb{R}$ be a Lebesgue measurable function such that
\[ f(x+y) = f(x) + f(y) \]
Show that $f(x) = ax$ for some $a \in \mathbb{R}$. \textbf{Hint:} Apply Lusin's theorem and the previous problem.
\end{question}
\begin{proof}
Use Lusin's theorem to get a closed set $F$ with $m(F^c) < 1/2$ and a continuous $g: F \to \R$ such that $f = g$ on $F$. By question 1, we can also restrict $F$ to a compact set while keeping $m(F) > 0$. We can then apply question 4 to get an open interval $(-M, M)$ contained in $F - F$. Since $F$ is compact, $g$ is also uniformly continuous on $F$, so given an $\varepsilon > 0$, there is some $\delta > 0$ for which $|g(x) - g(y)| < \varepsilon$ for any $x, y \in F, |x - y| < \delta$. For any $h \in (-\delta, \delta)$, we take $x \in F$ for which $h = (x + h) - x \in F - F$ and get $f(h) = f(x + h) - f(x) = g(x + h) - g(x) \implies |f(h)| < \varepsilon$. This shows that $f$ is continuous at 0. For any arbitrary $x_0 \in \R$, we can write $f(x + x_0) = f(x) + f(x_0)$ which shows that $f$ is continuous at $x_0$, therefore $f$ is continuous on $\R$. By continuity, it then suffices to only show that $f(x) = ax$ on $x \in \Q$.

Set $a := f(1)$. $f(0) = f(0) + f(0) \implies f(0) = 0$ and for any $x \in \Q$, we have $f(-x) = f(0) - f(x) = -f(x)$. For any $n \in \N$, we have $f(n) = \sum_{i = 1}^n f(1) = an$, and $f(-n) = -f(n) = -an$, so $f(x) = ax$ on $\Z$. For any $n \in \N$, we have $nf(1/n) = f(1) = a \implies f(1/n) = a/n$ and for any $p, q \in \N$, we have $qf(p/q) = f(p) = ap \implies f(p/q) = ap/q$, so for any $x \in \Q, f(x) = ax$. As noted earlier, continuity of $f$ shows that $f(x) = ax$ on $\R$.
\end{proof}

\begin{question}
Construct a non-linear $f : \mathbb{R} \to \mathbb{R}$ obeying $f(x+y) = f(x) + f(y)$ for all $x, y \in \mathbb{R}$. (Hint: consider $\mathbb{R}$ as a vector space with scalar field $\mathbb{Q}$)
\end{question}
\begin{proof}
Take a basis $\beta$ for $\R$ as a vector space over $\Q$. This gives us a unique expression of any real number $x$ as $q_1e_1 + \cdots + q_ne_n$ for $q_i \in \Q \setminus \{0\}, e_i \in \beta$. This lets us define $f$ by $q_1e_1 + \cdots + q_ne_n \mapsto q_1 + \cdots + q_n$. This function is obviously additive, but it cannot be continuous on $\R$, since its image is the disconnected set $\Q$. Hence it is not linear.
\end{proof}

\section{Royden Chapter 3}

\begin{question}[Royden 3.4]
Suppose $f$ is a real-valued function on $\mathbb{R}$ such that $f^{-1}(c)$ is measurable for each number $c$. Is $f$ necessarily measurable?
\end{question}
\begin{proof}
No. Let $E$ be any non-measurable set and define $f(x) = e^x(2\chi_E(x) - 1)$. Notice that $f^{-1}((0, \infty)) = E$ and $f^{-1}((-\infty, 0)) = E^c$, so we already see that $f$ is not measurable. Furthermore, if $f(x) = f(y)$, we have either $x, y \in E$ or $x, y \in E^c$, so in either case $e^x = e^y \implies x = y$, and hence $f^{-1}(c)$ for any $c \in \R$ is either empty or a singleton, hence measurable.
\end{proof}

\begin{question}[Royden 3.5]
Suppose the function $f$ is defined on a measurable set $E$ and has the property that $\{x \in E \mid f(x) > c\}$ is measurable for each rational number $c$. Is $f$ necessarily measurable?
\end{question}
\begin{proof}
It suffices to show that $f^{-1}((c, \infty))$ is measurable for any real $c$. By density of $\Q$, we can take a rational sequence $c_n > c$ such that $c_n \to c$. We then see that $\bigcup_{n = 1}^\infty (c_n, \infty) = (c, \infty)$ and hence $f^{-1}((c, \infty)) = \bigcup_{n = 1}^\infty f^{-1}((c_n, \infty))$ is measurable as a countable union of measurable sets.
\end{proof}

\begin{question}[Royden 3.6]
Let $f$ be a function with measurable domain $D$. Show that $f$ is measurable if and only if the function $g$ defined on $\mathbb{R}$ by 
\[
g(x) = 
\begin{cases}
f(x) & \text{for } x \in D \\
0 & \text{for } x \notin D
\end{cases}
\]
is measurable.
\end{question}
\begin{proof}
Note that for $c \geq 0$, we have $g^{-1}((c, \infty)) = f^{-1}(c, \infty)$ and for $c < 0$, $g^{-1}((c, \infty)) = D^c \cup f^{-1}((c, \infty)) \\\iff f^{-1}((c, \infty)) = g^{-1}((c, \infty)) \cap D$. The result follows after observing that $D, D^c$ are measurable.
\end{proof}

\begin{question}[Royden 3.13]
A real-valued measurable function is said to be \textit{semisimple} provided it takes only a countable number of values. Let $f$ be any measurable function on $E$. Show that there is a sequence of semisimple functions $\{f_n\}$ on $E$ that converges to $f$ uniformly on $E$.
\end{question}
\begin{proof}
We note that if $f$ is extended real-valued then this is not possible since $|f(x) - \varphi(x)| = \infty$ for any $x$ such that $f(x) = \pm \infty$ and any real-valued $\varphi$. We thus assume that $f$ is real-valued. Also by writing $f$ as a difference of its positive and negative parts (which are measurable), we may assume that $f$ is nonnegative.

Define $E_n$ to be the measurable set $E_n := f^{-1}([n, n + 1))$, which are mutually disjoint and cover $E$. For each $n$, the restriction $f_n$ of $f$ to $E_n$ is a bounded measurable function, so for any $\varepsilon > 0$, $f$ can be approximated by simple functions $\varphi_n^\varepsilon, \psi_n^\varepsilon$ on $E_n$, that is $\varphi_n^\varepsilon \leq f_n \leq \psi_n^\varepsilon$ and $0 \leq \psi_n^\varepsilon - \varphi_n^\varepsilon < \varepsilon$. Also by extending $\varphi_n^\varepsilon, \psi_n^\varepsilon$ to $E$ in the way of Royden 3.6, we get that $\varphi_n^\varepsilon, \psi_n^\varepsilon$ are simple functions on $E$. We define the semisimple functions $\varphi^\varepsilon = \sum_{i = 0}^\infty \varphi_n^\varepsilon, \psi^\varepsilon = \sum_{i = 0}^\infty \psi_n^\varepsilon$ and note that their restrictions to each $E_n$ are precisely $\varphi_n^\varepsilon, \psi_n^\varepsilon$. This implies that $\varphi^\varepsilon \leq f \leq \psi^\varepsilon$ and $0 \leq \psi^\varepsilon - \varphi^\varepsilon < \varepsilon$ on $E$, and hence $0 \leq f - \varphi^\varepsilon \leq \psi^\varepsilon - \varphi^\varepsilon < \varepsilon$. The result follows by taking the sequence $\varphi^{1/n}$.
\end{proof}

\begin{question}[Royden 3.14]
Let $f$ be a measurable function on $E$ that is finite a.e. on $E$ and $m(E) < \infty$. For each $\epsilon > 0$, show that there is a measurable set $F$ contained in $E$ such that $f$ is bounded on $F$ and $m(E \setminus F) < \epsilon$.
\end{question}
\begin{proof}

\end{proof}

\begin{question}[Royden 3.26]
For the function $f$ and the set $F$ in the statement of Lusin's Theorem, show that the restriction of $f$ to $F$ is a continuous function. Must there be any points at which $f$, considered as a function on $E$, is continuous?
\end{question}
\begin{proof}

\end{proof}

\begin{question}[Royden 3.27]
Show that the conclusion of Egoroff's Theorem can fail if we drop the assumption that the domain has finite measure.
\end{question}
\begin{proof}

\end{proof}

\begin{question}[Royden 3.29]
Prove the extension of Lusin's Theorem to the case that $E$ has infinite measure.
\end{question}
\begin{proof}

\end{proof}

\end{document}